\clearpage
\fancyhead[L]{\small ZACH. TABLES D'ABERRATION.}
\hypertarget{native-551-lower}{}\label{native-551-lower}
\begin{center}\small Göttingische gelehrte Anzeigen. No. 196. pp. 1956..1960. December 7, 1812.\\[3pt]\rule{.72\linewidth}{.3pt}\end{center}

\emph{Nouvelles tables d'aberration et de nutation pour quatorze cent quatre étoiles, avec une table générale d'aberration pour les planètes et les comètes, précédées d'une instruction qui renferme l'explication de l'usage de ces tables, suivies de plusieurs nouvelles tables destinées à faciliter les calculs astronomiques, par le Baron} \textsc{de Zach}. De l'imprimerie de Madame Mine et Comp. Marseille. 52 and 136 pages in octavo.

Since the discovery of aberration and nutation, astronomers have sought in many ways to ease the tiresome calculation of these small inequalities that confronts the observing astronomer daily. Various general tables have been arranged for this purpose, with varying degrees of success. There are also special tables for individual stars that give the aberration and nutation directly from the positions of the Sun and the lunar node at the time. The tables \SourceLeaf{552}of the latter kind by \textsc{Mezger}, \textsc{von Zach} and \textsc{Cagnoli} are in the hands of all practicing astronomers. They still cover so few stars, however, that one usually has to fall back on the general tables. In the present work Mr. \textsc{von Zach} has sought to make the calculations for a considerably larger number of stars somewhat easier than even the best-arranged general tables allow. Special tables in the strict sense for so many stars would be far too bulky and expensive, and he has done without them. The idea is very simple and has already been used by other astronomers in connection with aberration. The aberration in right ascension of any star can easily be put in the form $m\sin(\odot+A)$, where $m$ denotes the star's greatest positive aberration, $\odot$ the Sun's longitude at the time, and $A$ the complement to $360^\circ$ of the particular solar longitude at which the aberration vanishes and passes from a negative to a positive value. The quantities $m$ and $A$ depend on the mean position of the star. Since very great accuracy is unnecessary for them, once their values have been determined they can be used unchanged for quite a number of years. The same holds for aberration in declination, nutation in right ascension, and nutation in declination, except that $m$ and $A$ naturally have different values for each of these elements. Mr. \textsc{von Zach} gives these four values of $A$ and the logarithms of the corresponding $m$, the former to minutes and the latter to four decimal places, for 1404 of the most prominent stars from all the constellations visible in Europe, using their mean positions for 1800. This takes 94 pages and reduces the whole task to looking up four logarithms of sines and four numbers obtained from logarithms. One thus has to consult the logarithmic tables in four places, one more than when using the \emph{general} tables of \textsc{Gauss}, but has less to write altogether. For every star included in the new tables of \textsc{von Zach}, one will therefore gladly prefer them to the general ones. In Mr. \textsc{von Zach}'s otherwise very instructive introduction, a slight error has crept into the theoretical derivation, where the auxiliary quantity $A$ is assigned a wrong meaning. --- After the tables for 1404 stars, the 36 so-called \textsc{Maskelyne} stars are listed again. First come their right ascensions according to \textsc{Maskelyne}, together with $A$ and $\log m$ in time for aberration and nutation. Then come the declinations of the same stars according to \textsc{Piazzi}, together with the two corresponding values of $A$ and $\log m$, all for 1802. We note here some differences from the first table, admittedly insignificant in themselves, that cannot be explained by the change of epoch alone and seem to be calculation errors. For example, for the declination of $\alpha$ Tauri we have \[ \begin{array}{r@{}l|c|c} \multicolumn{2}{c|}{} & \text{first table} & \text{second table}\\ \text{Aberration}\; & \left\{\begin{array}{l}\log m\\ A\end{array}\right. & \begin{array}{c}0\,5806\\7^{\mathrm s}\;23^\circ\;26'\end{array} & \begin{array}{c}0.5772\\7^{\mathrm s}\;23^\circ\;15'\end{array}\\ \text{Nutation}\; & \left\{\begin{array}{l}\log m\\ A\end{array}\right. & \begin{array}{c}0.9673\\3^{\mathrm s}\;18^\circ\;18'\end{array} & \begin{array}{c}0.9682\\3^{\mathrm s}\;18^\circ\;6'\end{array} \end{array} \]

A similar table follows for the stars $\alpha$ and $\beta$ in Ursa Minor, based on their positions in 1790, 1800, 1810 and 1820.

Although the quantities $m$ and $A$ strictly apply only to the year for which they were calculated, they can still be kept for quite a long stretch of years before and after that epoch without introducing a significant error. The author gives a couple of special tables for estimating how much right ascension and declination can change before the resulting error reaches 1 second. This table also helps one judge how far \SourceLeaf{553}one may use the constants of a neighboring star from the large table for a star outside the 1404 that it lists. We pass over without comment a few following tables for aberration, nutation, precession and correction of noon from corresponding solar altitudes, since the author has already given them elsewhere and they are familiar. --- Table XX makes it easy to find the Sun's longitude as an argument for aberration for a day in a year for which one has no astronomical almanac at hand. It shows by how much the epoch of the mean solar longitude in each year from 1700 to 1827 is greater than in 1803. With sufficient accuracy for the stated purpose, one can also regard these numbers as the differences of the true longitudes on days with the same calendar date. One can thus readily derive the solar longitudes for any of those years from those of any other year for which one happens to have an astronomical almanac at hand. Likewise, Mr. \textsc{von Zach} gives a table of the longitude of the lunar node at the beginning of every year from 1700 to 1827 so that the ephemerides for any one of those years can be transferred to another. Last comes a table for converting sidereal time to mean solar time, arranged somewhat more conveniently than the author's familiar small tables calculated for the same purpose, which are in the hands of all astronomers. There is also a small table for the part of nutation that depends on the Sun's longitude. That part could indeed have been combined with Table XXII at once, which would have shortened the calculation a little further. The author presumably left it out deliberately so that this table would give the uncorrected mean right ascensions of the Sun. For the same reason, it seems, he omitted the addition of a negative constant, which is otherwise used to make the lunar nutation always positive.

This brief account of the contents shows that with this collection of tables the distinguished author has once again earned the thanks of all practicing astronomers, who must greatly welcome anything that makes their laborious daily calculations easier.
